\section*{Capítulo \ref{ch:g12:past-climates} --- Leer los climas del pasado}

\begin{solution}{exo:g12:past-climates:1}
(a) un \omterm{def:g12:past-climates:proxy}{indicador indirecto}; (b) un \omterm{def:g12:past-climates:ice-core}{testigo de hielo}; (c) un \omterm{def:g12:past-climates:glacial-cycle}{interglacial}; (d) el \omterm{def:g12:past-climates:pollen}{análisis polínico};
(e) el \omterm{def:g12:past-climates:lgm}{Último Máximo Glacial}.
\end{solution}

\begin{solution}{exo:g12:past-climates:2}
(a) El hielo polar: el dióxido de carbono y el metano del aire de la época.
(b) Los sedimentos lacustres o la turba: la vegetación, aquí árboles que necesitan veranos
cálidos. (c) Los sedimentos oceánicos: el volumen de hielo sobre los continentes y la temperatura
del \omterm{def:g10:layered-ocean:thermocline}{océano profundo}. (d) Los árboles: las condiciones de crecimiento, año a año.
\end{solution}

\begin{solution}{exo:g12:past-climates:3}
(a) $-50 - (-120) = \qty{70}{m}$ en \num{7000} años: $\qty{70000}{mm}/\num{7000}
= \qty{10}{mm}$ al año. (b) $10/3.7 \approx 2.7$: la subida media de la
deglaciación fue unas 2.7 veces más rápida que la de 2006--2018.
\end{solution}

\begin{solution}{exo:g12:past-climates:4}
Un siglo son diez décadas: $\qty{2.0}{\degreeCelsius}$ por siglo. Mil
años son cien décadas: $\qty{20}{\degreeCelsius}$ por cada mil años.
\end{solution}

\begin{solution}{exo:g12:past-climates:5}
(a) \num{23000} años: la precesión del eje; \num{42000} años: la
inclinación del eje; unos \num{100000} años: la forma de la órbita.
(b) $\num{1000000}/\num{41000} \approx 24$ ciclos; $\num{1000000}/\num{100000}
= 10$ ciclos.
\end{solution}

\begin{solution}{exo:g12:past-climates:6}
(a) $266.6 + 7.83 \times (-9.4) \approx 266.6 - 73.6 = \qty{193}{ppm}$: solo
\qty{3}{ppm} por encima de los \qty{190}{ppm} medidos. (b) $266.6 + 7.83 \times 5.5
\approx \qty{310}{ppm}$, \qty{11}{ppm} por encima del valor medido más alto: el
punto más cálido queda por debajo de la recta. (c) $\Delta T = (427 - 266.6)/7.83 \approx
\qty{20}{\degreeCelsius}$. La recta solo describe cómo variaron juntas las dos magnitudes,
en ciclos naturales, dentro del intervalo de los datos
(de \qty{-11}{\degreeCelsius} a \qty{+6}{\degreeCelsius}, de 174 a
\qty{299}{ppm}); no es una ley, y un valor muy alejado de ese intervalo no puede
leerse en ella.
\end{solution}

\begin{solution}{exo:g12:past-climates:7}
(a) Hierbas, \qty{84}{\percent}: un paisaje abierto casi sin árboles, frío.
(b) Roble (\qty{42}{\percent}) y avellano (\qty{34}{\percent}): un bosque de
frondosas, con veranos cálidos. (c) El abedul se extiende primero, luego el pino; el avellano y el roble
llegan los últimos. (d) Los árboles solo se extienden tan deprisa como viajan sus semillas y crecen sus
árboles jóvenes, a partir de los lugares donde sobrevivieron al frío; además, tiene
que formarse un suelo. El bosque puede ir con retraso respecto al clima.
\end{solution}

\begin{solution}{exo:g12:past-climates:8}
(a) \qty{300}{cm}, es decir, \qty{3000}{mm}, en \num{8000} años: $3000/8000 =
\qty{0.375}{mm}$ al año. (b) $\num{6000} + (470 - 300)/300 \times \num{8000}
\approx \num{6000} + \num{4530} \approx \num{10500}$ años. (c) $10/0.375
\approx 27$ años.
\end{solution}

\begin{solution}{exo:g12:past-climates:9}
(a) $7.16 - 4.11 = 3.05$ millones de kilómetros cuadrados, y $3.05/4.11 \approx
\qty{74}{\percent}$ más de tierra. (b) Sí: en el \omterm{def:g12:past-climates:lgm}{Último Máximo Glacial} el mar
estaba de 125 a \qty{134}{m} por debajo del actual, mucho más que los \qty{36}{m}
que unen Borneo con Asia en el mapa. (c) El mapa usa el fondo marino actual, que no ha
seguido exactamente igual desde entonces (se han depositado sedimentos sobre él, y el
terreno se ha movido un poco); la malla es además gruesa, un punto cada 0.1
grados, y el nivel más bajo del mar solo se conoce con una precisión de unos diez metros.
\end{solution}

\begin{solution}{exo:g12:past-climates:10}
(a) El valor del \omterm{def:g12:past-climates:lgm}{Último Máximo Glacial}, \qty{4.92}{\permil}, es el
más alto. Es lo esperado: los grandes inlandsis de la época almacenaban agua pobre
en oxígeno 18, lo que dejaba el océano, y las conchas crecidas en él, más ricos en
oxígeno 18; además, el agua profunda estaba más fría. (b) Para que las épocas frías y con mucho hielo queden
abajo y las cálidas arriba, como en una curva de temperatura.
\end{solution}

\begin{solution}{exo:g12:past-climates:11}
(a) Una correlación muestra que los dos registros suben y bajan a la vez, cerca
de una recta; no muestra cuál causó cuál, ni que alguno de ellos causara
al otro. (b) Un clima más cálido podría haber liberado dióxido de carbono al aire;
o una tercera causa, como los cambios de la órbita de la Tierra, podría haber movido
a ambos; o cada uno podría haber empujado al otro. (c) La cronología de los cambios,
es decir, cuál vino primero, a partir de una datación precisa del hielo y de su aire; y una
comprensión de cómo actúa el dióxido de carbono sobre la radiación de la Tierra
(el \cref{ch:g12:climate-change}).
\end{solution}

\begin{solution}{exo:g12:past-climates:12}
(a) $\num{1300000}/\num{41000} \approx 32$ ciclos entre hace 2.6 y 1.3
millones de años; $\num{700000}/\num{98500} \approx 7$ ciclos en los últimos
\num{700000} años. (b) \num{41000} años: la inclinación del eje; unos
\num{100000} años: la forma de la órbita. (c) Un ciclo de \num{41000} años
se muestrea unas $\num{41000}/\num{2000} \approx 20$ veces, lo bastante para dibujarlo;
un cambio que dura un siglo cae entre dos muestras, o dentro de una, y
casi no deja huella.
\end{solution}

\begin{solution}{exo:g12:past-climates:13}
(a) $192/10 \approx 19$ y $198/10 \approx 20$: unas 20 veces más rápido.
(b) Frente a \qty{4}{ppm} por siglo: de $192/4 = 48$ a $198/4 \approx 50$
veces; frente a \qty{42}{ppm}: de $192/42 \approx 4.6$ a $198/42 \approx 4.7$
veces. El aumento actual es de 5 a 50 veces más rápido que el del PETM.
(c) El aumento es de $2000 - 900 = \qty{1100}{ppm}$: $1100/42 \approx 26$ siglos,
unos \num{2600} años al ritmo más rápido del PETM; $1100/195 \approx 5.6$
siglos, menos de 600 años, al ritmo actual.
\end{solution}

\begin{solution}{exo:g12:past-climates:14}
(a) Diez veces de 1 a \qty{1.5}{\degreeCelsius}: de 10 a \qty{15}{\degreeCelsius}
por cada mil años, es decir, de 0.10 a \qty{0.15}{\degreeCelsius} por década.
(b) El ritmo global actual, \qty{0.20}{\degreeCelsius} por década, es de 1.3 a 2
veces el calentamiento más rápido sobre Groenlandia. (c) Un solo lugar puede cambiar mucho
más que la media del globo, donde los cambios de las distintas regiones se
compensan en parte: un ritmo global se compara con un ritmo global.
\end{solution}

\begin{solution}{exo:g12:past-climates:15}
(a) \qty{14}{m} a \qty{12}{m} por cada mil años llevan $14/12 \approx
1.17$ miles de años, unos \num{1170} años: hace unos \num{9800} años.
(b) No: el mar tenía que mantenerse al menos \qty{36}{m} por debajo del actual para unir
Borneo con Asia; en cuanto pasó de \qty{-36}{m}, el puente de tierra quedó sumergido, y
Borneo volvió a ser una isla. (c) $36/12 = 3$ miles de años, unos
\num{3000} años; en realidad, la subida se frenó a partir de hace unos \num{8200} años,
así que tardó más.
\end{solution}

\begin{solution}{pb:g12:past-climates:1}
\textbf{1.} El propio hielo, mediante sus isótopos pesados, registra la
temperatura; el aire de sus burbujas registra el dióxido de carbono.

\textbf{2.} $5.5 - (-10.6) = \qty{16.1}{\degreeCelsius}$.

\textbf{3.} $299 - 174 = \qty{125}{ppm}$.

\textbf{4.} Los puntos quedan cerca de una recta: los dos registros subieron
y bajaron a la vez. No dice cuál causó cuál.

\textbf{5.} Los registros del pasado se dibujan con las épocas más antiguas a uno u otro
lado; leído al revés, un calentamiento se tomaría por un enfriamiento.

\textbf{6.} Un \omterm{def:g12:past-climates:glacial-cycle}{periodo glacial} es un intervalo frío dentro de una glaciación, en el que los
inlandsis crecen mucho más allá de su tamaño actual; un \omterm{def:g12:past-climates:glacial-cycle}{interglacial} es un intervalo
cálido entre dos \omterm{def:g12:past-climates:glacial-cycle}{periodos glaciales}, con inlandsis parecidos a los de hoy o
más pequeños.

\textbf{7.} En un \omterm{def:g12:past-climates:glacial-cycle}{interglacial}, que empezó hace unos \num{11000} años.

\textbf{8.} $\num{800000}/\num{98000} \approx 8$: unos ocho ciclos.

\textbf{9.} La forma de la órbita.

\textbf{10.} Un ritmo de \num{41000} años, el periodo de la inclinación del
eje.

\textbf{11.} $-0.4 - (-9.4) = \qty{9.0}{\degreeCelsius}$.

\textbf{12.} $264 - 190 = \qty{74}{ppm}$.

\textbf{13.} \qty{6}{\degreeCelsius}; el cambio antártico, \qty{9.0}{\degreeCelsius},
es mayor, unas 1.5 veces el global.

\textbf{14.} $\num{18000} - \num{11000} = \num{7000}$ años.

\textbf{15.} $6/7 \approx \qty{0.86}{\degreeCelsius}$ por cada mil años.

\textbf{16.} Mil años son cien décadas: $0.86/100 \approx
\qty{0.0086}{\degreeCelsius}$ por década.

\textbf{17.} $0.20 \times 100 = \qty{20}{\degreeCelsius}$ por cada mil años.

\textbf{18.} $20/0.86 \approx 23$ veces.

\textbf{19.} $20/1.5 \approx 13$ y $20/1 = 20$: de 13 a 20 veces.

\textbf{20.} El calentamiento global actual es \textbf{unas veinte veces más rápido} que
el calentamiento a la salida de la última glaciación: 23 veces su ritmo medio, de 13 a 20 veces
su ritmo más rápido.
\end{solution}
